\section{One value for each retained measurement}
\label{sec:joint-source-measurements}

The current-assets total appears in more than one balance-sheet relation.
Its components must add to it, and it must combine with noncurrent assets
to give total assets. If these relations are tested separately, each test
can choose a different unrounded current-assets amount. A company has one
amount at that date, so the question is whether the same choice can satisfy
all the selected relations. This section supplies that narrower joint test
before attempting to construct an accounting history.

Let \(d_i\) denote a retained displayed amount and \(x_i\) its proposed
underlying amount, both expressed in millions of US dollars. The index
identifies an observation, including its company, perimeter, reporting basis
and measurement date. Two uses of that observation share \(x_i\); a
different date or source identity does not acquire the same variable merely
because its caption is similar. The retained balance example concerns
Apple's consolidated amounts at 28 September and 28 December 2024. It
preserves the two dates and does not equate their balances.

To constrain an underlying amount from its display, the example proposes
rounding to the nearest million, with exact half ties rounded away from
zero. This is the named hypothesis introduced for the partial-state example,
not a rounding rule established from the issuer's disclosures. In displayed
million-dollar units, it implies
\begin{equation}
\label{eq:joint-source-signed-rounding}
 I(d_i)=
 \begin{cases}
 [d_i-\tfrac12,\,d_i+\tfrac12), & d_i>0,\\
 (d_i-\tfrac12,\,d_i+\tfrac12], & d_i<0,\\
 (-\tfrac12,\,\tfrac12), & d_i=0.
 \end{cases}
\end{equation}
For a positive display, its lower half tie rounds upward to that display,
whereas its upper half tie rounds to the next larger integer. The endpoints
reverse for a negative display. Neither half tie rounds to zero, which
explains the open interval in the third line. If the hypothesis is absent,
the source adapter supplies no interval and the joint question remains
unresolved. A literal dash also remains nonnumeric; it is not converted into
the third line.

\subsection{A common assignment must satisfy every selected equation}

Each selected relation supplies exact signed coefficients: components occur
on the left and the reported target is moved to the left with a negative
coefficient. A separate named hypothesis asserts that this displayed
relation also describes the underlying unrounded quantities. With one row
per selected relation and one column per distinct observation, these
coefficients form a rational matrix \(A\). The conditional question is
\begin{equation}
\label{eq:joint-source-feasible-set}
 \mathcal F(A,I)=
 \left\{x\in\mathbb R^n:
 x_i\in I(d_i)\ \text{for every }i,\quad Ax=0\right\}.
\end{equation}
An element of \(\mathcal F(A,I)\) uses the same value everywhere a
measurement recurs. This is the difference from checking whether each
individual residual interval contains zero. For example, the two equalities
\(x=0\) and \(x=1\) are separately compatible with \(x\in[0,1]\),
but jointly require \(0=1\). Their intersection is empty even though both
local tests pass.

In the retained annual current-assets relation, the displayed components are
cash and cash equivalents 29,943; current marketable securities 35,228;
accounts receivable 33,410; vendor non-trade receivables 32,833; inventories
7,286; and other current assets 14,287. Their sum is
\begin{equation}
\label{eq:joint-source-current-assets}
 29{,}943+35{,}228+33{,}410+32{,}833+7{,}286+14{,}287
 =152{,}987.
\end{equation}
Thus the displayed values themselves satisfy this selected equation and lie
inside their proposed rounding intervals. The joint calculation also uses
the same current-assets variable in the total-assets equation, and the same
total-assets variable in the balance-sheet identity. It tests all these
conditions together rather than supplying a new value to each equation.

The reference implementation \path{JointMeasurementConstraints/1} accepts
explicit rational coefficients and intervals, including endpoint inclusion.
It uses the existing exact solver to seek an assignment and checks each
returned rational value independently with exact fraction arithmetic.
For every variable, that check applies the original interval inequalities;
for every equation, it substitutes all values into the original terms and
requires exact equality. A valid witness establishes nonemptiness of this
declared set. When the solver reports infeasibility, its conflicting
constraints are checked again by the same solver. That recheck is useful
consistency evidence, but it is not an independently verified formal proof.

The supplied computational budget is part of the request. An exhausted
budget or an inconclusive solver response gives an unknown result, not
infeasibility. Missing premises are reported before solving, after the
structure and types have been checked. This prevents an unresolved
assumption from concealing a malformed equation or a unit mismatch.

\subsection{Preserve the observations when constructing the system}

The finite source adapter \path{RetainedJointSourceRelations/1} receives
the exact combined statement request, an explicit selection of supported
relations and optional hypotheses. It recomputes the preserved statement
description. Coefficients, observation identities, dates and units come from
that description; the caller cannot replace them with a more convenient
matrix or a fabricated source result. The currently supported selection is
the nine current balance relations at each of the two retained dates.

The observation's retained digest identifies its shared variable. Its
consolidated perimeter, reporting basis, original period description,
measurement kind and exact instant remain attached to its type. The source
unit \emph{USD million} is represented as a US-dollar monetary dimension
with scale \(10^6\). Coefficients and numerical interval endpoints remain
in displayed million-dollar units, so the conversion does not multiply
the amounts a second time. The complete original context is preserved
alongside this typed representation.

All 256 observations and all 143 original relation dispositions remain in
the returned description, including relations outside this first joint
selection. The six derived subtotals still lack independent observed targets,
and the two calculations that depend on unresolved dashes remain unresolved.
Cross-period comparisons, omitted classes of transaction and explanatory
differences retain their original meanings. None becomes an asserted
underlying identity simply because another part of the balance sheet has
a feasible assignment.

The earlier displayed-statement report therefore keeps its own annotation
that joint underlying feasibility is uncomputed. A separate result identifies
the selected conditional system, its hypotheses and its assessment. This
separation lets a reader recover what the source description originally
said while examining the additional mathematical question. An empty selection
reports that no system has been asserted; an unasserted rounding or identity
hypothesis leaves the selected question unresolved.

\subsection{What the retained balance example establishes}

With all 18 supported balance relations selected, the system contains
54 distinct monetary observations, 27 at each date. Under the explicit
rounding and identity hypotheses, it has an exact rational witness.
Independent checks substitute that witness directly into the retained source
terms and separately derived rounding bounds. The displayed values themselves
also provide a checked witness for this selected example. These two checks
establish compatibility of the selected measurement equations under the
proposed assumptions; they do not identify the issuer's unrounded amounts.

The local interval calculations are retained alongside the joint result
using exactly the same premises. All 18 local residual intervals contain
zero, and the joint witness resolves the further question of whether the
same amounts can work throughout the selected network. The constructed
\(x=0, x=1\) example remains necessary: agreement between the two kinds
of check in this company example does not make the local test sufficient
in general.

Source changes, altered coefficients or dates, renamed hypotheses and
changed hypothesis meanings are rejected. Result verification recomputes
the deterministic source mapping and checks the mathematical evidence.
It permits another valid witness because a feasible system can have many
solutions. It does not require the solver to choose the same assignment,
conflict ordering or elapsed time on every replay. Recorded checks also
exercise missing hypotheses, empty selections and zero computational budgets.
The executable definitions and complete case results are retained under
\path{data/accounting-master/}:
\begin{quote}\small
\path{joint-measurement-constraints-20261024}\\
\path{joint-source-relation-binding-20261024}
\end{quote}

\subsection{From compatible measurements to an executable company}

A feasible balance-sheet assignment still leaves the simulator's main
historical problem unresolved. Consolidated cash does not identify the
available cash of a particular legal book; an inventory caption does not
supply quantities, cohorts or cost layers; and aggregate liabilities do
not specify contracts, maturities or payment conditions. Those distinctions
affect which events can occur and how they change future statements.

The selected equations constrain measurements at two dates without supplying
the events between them. They assign neither probabilities to possible
histories nor behavioral rules to future transactions. Section~\ref{sec:temporal-statement-network} first connects these measurements
to dated stock-flow relations and a wider selected statement network.
Reconstruction must then retain those constraints while representing the
missing state and admissible event histories explicitly. Until that step
has its own source, accounting and computational support, the actual history
remains unknown and no executable company opening is produced.
